\documentclass[11pt]{article}
\usepackage{amsmath, amssymb, geometry, amsthm, graphicx, enumitem}
\usepackage[colorlinks=true, linkcolor=blue, urlcolor=blue]{hyperref}

\geometry{letterpaper, margin=1in}

\title{Homework 2.2: Separable Differential Equations}
\author{Tashfeen Omran}
\date{September 2026}

\begin{document}

\maketitle
\tableofcontents
\newpage

\section{Part 1: Comprehensive Introduction, Context, and Prerequisites}

\subsection{Core Concepts}
A differential equation is an equation that relates an unknown function to its derivatives. We use differential equations to model relationships where we know how a quantity changes (its rate of change) but not necessarily the explicit formula for the quantity itself. 

In this chapter, we focus on \textbf{First-Order Ordinary Differential Equations (ODEs)}, which involve only the first derivative, $\frac{dy}{dx}$ (or $y'$). A special, highly solvable class of these equations are known as \textbf{Separable Equations}. 

A first-order ODE is called \textit{separable} if it can be written in the form:
\[ \frac{dy}{dx} = g(x) \cdot p(y) \]
In other words, the derivative $\frac{dy}{dx}$ can be expressed as the product of a function that depends strictly on $x$ and a function that depends strictly on $y$. The name ``separable'' comes from the fact that we can algebraically manipulate the equation to ``separate'' the variables, moving all $y$-terms (including $dy$) to one side of the equation and all $x$-terms (including $dx$) to the other.

\subsection{Intuition and Derivation}
The intuition behind separable equations relies heavily on the chain rule and differentials. If we start with the separated form:
\[ \frac{1}{p(y)} dy = g(x) dx \]
Let $f(y) = \frac{1}{p(y)}$. We then integrate both sides with respect to their respective variables:
\[ \int f(y) dy = \int g(x) dx \]
Why does this work? Formally, if we have a function $y(x)$, the derivative of the anti-derivative of $f(y)$ with respect to $x$ is, by the chain rule, $f(y) \frac{dy}{dx}$. Integrating both sides with respect to $x$ gives $\int f(y) \frac{dy}{dx} dx = \int g(x) dx$, and the $dx$ algebraically cancels out to yield $\int f(y) dy$. This validates the incredibly useful trick of treating $\frac{dy}{dx}$ as if it were a fraction of differentials.

\subsection{Historical Context and Motivation}
The concept of separation of variables dates back to the late 17th century, developed shortly after Isaac Newton and Gottfried Wilhelm Leibniz invented calculus. Leibniz introduced the notation $\frac{dy}{dx}$ and differentials in 1675, and Johann Bernoulli formalized the method of separation of variables in 1694. The driving motivation was to solve physical problems related to motion, gravity, and curves. Early physicists needed a way to translate continuous rates of change (like a cooling object whose temperature drops at a rate proportional to the temperature difference) into explicit formulas for prediction. Separable equations were the very first class of differential equations to be systematically solved.

\subsection{Key Formulas}
\begin{itemize}
    \item \textbf{Separable Form:} $\frac{dy}{dx} = g(x) p(y)$
    \item \textbf{Separated Differentials:} $\frac{1}{p(y)} dy = g(x) dx$
    \item \textbf{General Solution Form:} $\int \frac{1}{p(y)} dy = \int g(x) dx + C$ 
\end{itemize}

\subsection{Prerequisites}
To master this topic, you must be entirely comfortable with:
\begin{enumerate}
    \item \textbf{Algebraic Factoring:} Knowing how to factor out variables, particularly recognizing exponent rules like $e^{a+b} = e^a e^b$.
    \item \textbf{Integration Techniques:} Specifically, basic power rule integration, $u$-substitution, and integration by parts.
    \item \textbf{Logarithms and Exponents:} A solid grasp of the properties of logarithms and exponentials (e.g., $e^{\ln(x)} = x$) is required, as solving for $y$ often involves exponentiating both sides.
\end{enumerate}

\newpage
\section{Part 2: Detailed Homework Solutions}

\subsection*{Problem 1}
\textbf{Determine whether the given differential equation is separable.}
\[ \frac{dy}{dx} + 10 \tan(x-y) = 0 \]

\textbf{Solution:}
First, we isolate the derivative:
\[ \frac{dy}{dx} = -10 \tan(x-y) \]
To be separable, we must be able to express the right-hand side in the form $g(x)p(y)$. 
Using trigonometric identities, there is no way to factor $\tan(x-y)$ into a product of a function solely of $x$ and a function solely of $y$. Because the $x$ and $y$ variables are locked inside the argument of the tangent function as a difference, they cannot be uncoupled via multiplication.
\newline
\textbf{Final Answer:} \textbf{D. No}

\vspace{0.5cm}
\subsection*{Problem 2}
\textbf{Determine whether the given differential equation is separable.}
\[ \frac{dy}{dx} = \frac{xe^{x+2y}}{x^2+2} \]

\textbf{Solution:}
We utilize the property of exponents, $e^{a+b} = e^a e^b$, to rewrite the numerator:
\[ \frac{dy}{dx} = \frac{x e^x e^{2y}}{x^2+2} \]
Now, we can clearly separate the terms involving $x$ from the terms involving $y$ by grouping:
\[ \frac{dy}{dx} = \left( \frac{xe^x}{x^2+2} \right) \cdot (e^{2y}) \]
This perfectly matches the form $\frac{dy}{dx} = g(x)p(y)$, where $g(x) = \frac{xe^x}{x^2+2}$ and $p(y) = e^{2y}$.
\newline
\textbf{Final Answer:} \textbf{B. Yes; because $\frac{dy}{dx} = g(x)p(y)$ where $g(x) = \frac{xe^x}{x^2+2}$ and $p(y) = e^{2y}$.}

\vspace{0.5cm}
\subsection*{Problem 3}
\textbf{Determine whether the given differential equation is separable.}
\[ (xy^4 + 3y^4)dy - 8x dx = 0 \]

\textbf{Solution:}
First, let's algebraically manipulate the equation to isolate $\frac{dy}{dx}$. 
Add $8x dx$ to both sides:
\[ (xy^4 + 3y^4)dy = 8x dx \]
Factor out the common term $y^4$ on the left side:
\[ y^4(x+3)dy = 8x dx \]
Divide both sides by $dx$ and by $y^4(x+3)$ to isolate $\frac{dy}{dx}$:
\[ \frac{dy}{dx} = \frac{8x}{y^4(x+3)} \]
We can write this as a product of two distinct functions:
\[ \frac{dy}{dx} = \left( \frac{8x}{x+3} \right) \cdot \left( \frac{1}{y^4} \right) \]
This matches the separable form $g(x)p(y)$ with $g(x) = \frac{8x}{x+3}$ and $p(y) = \frac{1}{y^4}$.
\newline
\textbf{Final Answer:} \textbf{B. Yes; because $\frac{dy}{dx} = g(x)p(y)$ where $g(x) = \frac{8x}{x+3}$ and $p(y) = \frac{1}{y^4}$.}

\vspace{0.5cm}
\subsection*{Problem 4}
\textbf{Solve the equation.}
\[ \frac{dx}{dt} = 4xt^3 \]
\textit{An implicit solution in the form $F(t,x)=C$ is requested.}

\textbf{Solution:}
Separate the variables $x$ and $t$. Divide by $x$ and multiply by $dt$:
\[ \frac{1}{x} dx = 4t^3 dt \]
Integrate both sides:
\[ \int \frac{1}{x} dx = \int 4t^3 dt \]
\[ \ln|x| = \frac{4t^4}{4} + C \]
\[ \ln|x| = t^4 + C \]
To write this in the requested implicit form $F(t,x)=C$, we subtract $t^4$ from both sides:
\[ \ln|x| - t^4 = C \]
\textbf{Final Answer:} \textbf{$\ln|x| - t^4 = C$}

\vspace{0.5cm}
\subsection*{Problem 5}
\textbf{Solve the equation.}
\[ \frac{dx}{dt} = \frac{4}{x e^{t+2x}} \]
\textit{An implicit solution in the form $F(t,x)=C$ is requested.}

\textbf{Solution:}
First, expand the exponential term in the denominator:
\[ \frac{dx}{dt} = \frac{4}{x e^t e^{2x}} \]
To separate variables, multiply both sides by $x e^{2x}$ and by $dt$:
\[ x e^{2x} dx = \frac{4}{e^t} dt \]
Rewrite with negative exponents for easier integration:
\[ x e^{2x} dx = 4e^{-t} dt \]
Integrate both sides:
\[ \int x e^{2x} dx = \int 4e^{-t} dt \]
For the left side, we must use Integration by Parts. Let:
\[ u = x \implies du = dx \]
\[ dv = e^{2x} dx \implies v = \frac{1}{2}e^{2x} \]
Using the formula $\int u \, dv = uv - \int v \, du$:
\[ \int x e^{2x} dx = \frac{1}{2}x e^{2x} - \int \frac{1}{2}e^{2x} dx \]
\[ = \frac{1}{2}x e^{2x} - \frac{1}{4}e^{2x} \]
Now integrate the right side:
\[ \int 4e^{-t} dt = -4e^{-t} + C \]
Equate the two integrated sides:
\[ \frac{1}{2}x e^{2x} - \frac{1}{4}e^{2x} = -4e^{-t} + C \]
To put it in the required form $F(t,x)=C$, move all variable terms to the left:
\[ \frac{1}{2}x e^{2x} - \frac{1}{4}e^{2x} + 4e^{-t} = C \]
\textbf{Final Answer:} \textbf{$\frac{1}{2}xe^{2x} - \frac{1}{4}e^{2x} + 4e^{-t} = C$}

\vspace{0.5cm}
\subsection*{Problem 6}
\textbf{Solve the equation.}
\[ (x^3 + 6x^3y^4)dx + e^{x^4}y^3dy = 0 \]
\textit{An implicit solution in the form $F(x,y)=C$ is requested.}

\textbf{Solution:}
First, move the $dy$ term to the right side:
\[ (x^3 + 6x^3y^4)dx = -e^{x^4}y^3dy \]
Factor $x^3$ from the left side:
\[ x^3(1 + 6y^4)dx = -e^{x^4}y^3dy \]
Separate the variables. Divide by $e^{x^4}$ and by $(1+6y^4)$:
\[ \frac{x^3}{e^{x^4}} dx = \frac{-y^3}{1+6y^4} dy \]
Rewrite the left side with a negative exponent:
\[ x^3 e^{-x^4} dx = \frac{-y^3}{1+6y^4} dy \]
Integrate both sides:
\[ \int x^3 e^{-x^4} dx = \int \frac{-y^3}{1+6y^4} dy \]
\textit{Left side evaluation ($u$-substitution):}
Let $u = -x^4 \implies du = -4x^3 dx \implies -\frac{1}{4}du = x^3 dx$.
\[ \int x^3 e^{-x^4} dx = \int -\frac{1}{4} e^u du = -\frac{1}{4}e^u = -\frac{1}{4}e^{-x^4} \]
\textit{Right side evaluation ($w$-substitution):}
Let $w = 1+6y^4 \implies dw = 24y^3 dy \implies \frac{1}{24}dw = y^3 dy$.
\[ \int \frac{-y^3}{1+6y^4} dy = \int \frac{-1}{24} \frac{1}{w} dw = -\frac{1}{24}\ln|w| = -\frac{1}{24}\ln(1+6y^4) \]
(Note: Absolute value is dropped because $1+6y^4 > 0$ for all real $y$).
Equating the integrals gives:
\[ -\frac{1}{4}e^{-x^4} = -\frac{1}{24}\ln(1+6y^4) + C_1 \]
The requested form in the image is $C = \ln(1+6y^4) - 6e^{-x^4}$. Let us manipulate our equation to match this form. Multiply the entire equation by $24$:
\[ -6e^{-x^4} = -\ln(1+6y^4) + 24C_1 \]
Add $\ln(1+6y^4)$ to both sides, and let $C = 24C_1$:
\[ \ln(1+6y^4) - 6e^{-x^4} = C \]
\textbf{Final Answer:} \textbf{$\ln(1+6y^4) - 6e^{-x^4} = C$}

\vspace{0.5cm}
\subsection*{Problem 7}
\textbf{Solve the initial value problem.}
\[ \frac{1}{\theta} \frac{dy}{d\theta} = \frac{y \sin\theta}{y^4+1}, \quad y(\pi) = 1 \]

\textbf{Solution:}
Separate the variables. Multiply by $\theta \, d\theta$ and multiply by $\frac{y^4+1}{y}$:
\[ \frac{y^4+1}{y} dy = \theta \sin\theta \, d\theta \]
Simplify the left side by splitting the fraction:
\[ \left( y^3 + \frac{1}{y} \right) dy = \theta \sin\theta \, d\theta \]
Integrate both sides:
\[ \int \left( y^3 + \frac{1}{y} \right) dy = \int \theta \sin\theta \, d\theta \]
The left side integrates straightforwardly via the power rule and natural logarithm:
\[ \frac{y^4}{4} + \ln|y| \]
For the right side, use Integration by Parts. Let:
\[ u = \theta \implies du = d\theta \]
\[ dv = \sin\theta \, d\theta \implies v = -\cos\theta \]
\[ \int \theta \sin\theta \, d\theta = -\theta \cos\theta - \int -\cos\theta \, d\theta = -\theta \cos\theta + \sin\theta + C \]
Combining both sides yields the general solution:
\[ \frac{y^4}{4} + \ln|y| = -\theta \cos\theta + \sin\theta + C \]
Now apply the initial condition $y(\pi) = 1$ to find $C$:
\[ \frac{(1)^4}{4} + \ln(1) = -\pi \cos(\pi) + \sin(\pi) + C \]
\[ \frac{1}{4} + 0 = -\pi(-1) + 0 + C \]
\[ \frac{1}{4} = \pi + C \implies C = \frac{1}{4} - \pi \]
Substitute $C$ back into the general equation:
\[ \frac{y^4}{4} + \ln|y| = -\theta \cos\theta + \sin\theta + \frac{1}{4} - \pi \]
\textbf{Final Answer:} \textbf{$\frac{y^4}{4} + \ln(y) = -\theta \cos(\theta) + \sin(\theta) + \frac{1}{4} - \pi$}

\vspace{0.5cm}
\subsection*{Problem 8}
\textbf{Solve the initial value problem.}
\[ \frac{dy}{dt} = 8t \cos^2(y), \quad y(1) = \frac{7\pi}{4} \]

\textbf{Solution:}
Separate variables. Divide by $\cos^2(y)$ and multiply by $dt$:
\[ \frac{1}{\cos^2(y)} dy = 8t \, dt \]
Using the trigonometric identity $\frac{1}{\cos^2(y)} = \sec^2(y)$:
\[ \sec^2(y) dy = 8t \, dt \]
Integrate both sides:
\[ \int \sec^2(y) dy = \int 8t \, dt \]
\[ \tan(y) = 4t^2 + C \]
Apply the initial condition $y(1) = \frac{7\pi}{4}$:
\[ \tan\left(\frac{7\pi}{4}\right) = 4(1)^2 + C \]
Recall that $\frac{7\pi}{4}$ is in the fourth quadrant where tangent is negative, specifically $\tan(\frac{7\pi}{4}) = -1$:
\[ -1 = 4 + C \implies C = -5 \]
Substitute $C$ back into the general solution:
\[ \tan(y) = 4t^2 - 5 \]
\textbf{Final Answer:} \textbf{$\tan(y) = 4t^2 - 5$}

\vspace{0.5cm}
\subsection*{Problem 9}
\textbf{Solve the initial value problem.}
\[ \frac{dy}{dx} = x^5(y-4), \quad y(0) = 6 \]

\textbf{Solution:}
Separate variables by dividing by $(y-4)$ and multiplying by $dx$:
\[ \frac{1}{y-4} dy = x^5 dx \]
Integrate both sides:
\[ \int \frac{1}{y-4} dy = \int x^5 dx \]
\[ \ln|y-4| = \frac{x^6}{6} + C \]
To solve for $y$ explicitly, exponentiate both sides (use base $e$):
\[ |y-4| = e^{\frac{x^6}{6} + C} \]
\[ |y-4| = e^C \cdot e^{\frac{x^6}{6}} \]
Let $A = \pm e^C$. Thus, we can drop the absolute value:
\[ y - 4 = A e^{\frac{x^6}{6}} \]
\[ y = 4 + A e^{\frac{x^6}{6}} \]
Use the initial condition $y(0) = 6$ to find $A$:
\[ 6 = 4 + A e^{\frac{0^6}{6}} \]
\[ 6 = 4 + A(1) \implies A = 2 \]
Substitute $A$ back into the explicit solution:
\[ y = 4 + 2e^{\frac{x^6}{6}} \]
\textbf{Final Answer:} \textbf{$y - 4 = 2 e^{\frac{x^6}{6}}$ (or equivalently, $y = 4 + 2 e^{\frac{x^6}{6}}$)}

\vspace{0.5cm}
\subsection*{Problem 10}
\textbf{Solve the initial value problem.}
\[ \sqrt{y} dx + (2+x)dy = 0, \quad y(-1) = 49 \]

\textbf{Solution:}
First, rearrange the equation to isolate the differentials on opposite sides:
\[ (2+x)dy = -\sqrt{y} dx \]
Separate the variables. Divide by $\sqrt{y}$ and by $(2+x)$:
\[ \frac{1}{\sqrt{y}} dy = \frac{-1}{2+x} dx \]
Rewrite the radical as a fractional exponent:
\[ y^{-1/2} dy = -\frac{1}{2+x} dx \]
Integrate both sides:
\[ \int y^{-1/2} dy = \int -\frac{1}{2+x} dx \]
Apply the power rule to the left side and natural log to the right:
\[ \frac{y^{1/2}}{1/2} = -\ln|2+x| + C \]
\[ 2\sqrt{y} = -\ln|2+x| + C \]
Move the logarithm term to the left side to match the requested format:
\[ 2\sqrt{y} + \ln|2+x| = C \]
Apply the initial condition $y(-1) = 49$:
\[ 2\sqrt{49} + \ln|2 + (-1)| = C \]
\[ 2(7) + \ln|1| = C \]
\[ 14 + 0 = C \implies C = 14 \]
Substitute $C$ back into the equation:
\[ 2\sqrt{y} + \ln|2+x| = 14 \]
Since $x$ is near $-1$ (from the initial condition), $2+x > 0$, so we can write parentheses instead of absolute values.
\newline
\textbf{Final Answer:} \textbf{$2\sqrt{y} + \ln(2+x) = 14$}

\newpage
\section{Part 3: In-Depth Analysis of Problems and Techniques}

\subsection{A) Problem Types and General Approach}
Reviewing the homework, we can categorize the problems into three major distinct types:

\begin{enumerate}
    \item \textbf{Identification and Setup (Problems 1-3):}
    \begin{itemize}
        \item \textit{Approach:} Use algebraic manipulation (distributive property, exponent rules) to see if $f(x,y)$ can be factored completely into $g(x) \cdot p(y)$. If an $x$ and $y$ are tied together inside a function via addition or subtraction (like $\tan(x-y)$), they are inseparable.
    \end{itemize}
    \item \textbf{General Implicit Solutions (Problems 4-6):}
    \begin{itemize}
        \item \textit{Approach:} Separate variables, perform integration on both sides, and add $+C$ to one side (usually the $x$-side). Then, rather than solving for $y$ explicitly, manipulate the equation algebraically so all variables are on one side and only the constant $C$ remains on the right.
    \end{itemize}
    \item \textbf{Initial Value Problems (IVPs) (Problems 7-10):}
    \begin{itemize}
        \item \textit{Approach:} Solve the differential equation as usual to find the general solution. Immediately after integrating, plug in the given initial condition $(x_0, y_0)$ to solve for the specific constant $C$. Once $C$ is known, rewrite the final equation.
    \end{itemize}
\end{enumerate}

\subsection{B) Key Algebraic and Calculus Manipulations}
\begin{itemize}
    \item \textbf{Splitting Exponentials (Problem 2, 5):} Recognizing that $e^{x+2y} = e^x \cdot e^{2y}$ is mandatory. Without this step, the variables appear un-separable.
    \item \textbf{Factoring (Problem 3, 6):} Identifying common terms, such as factoring $y^4$ out of $xy^4 + 3y^4$ to get $y^4(x+3)$, separates the variables geometrically.
    \item \textbf{Integration by Parts (Problem 5, 7):} Crucial for integrals involving products of polynomials and transcendental functions, such as $\int x e^{2x} dx$ and $\int \theta \sin\theta \, d\theta$. Remember the LIATE rule to choose your $u$.
    \item \textbf{Trigonometric Identities (Problem 8):} The equation yields $1/\cos^2(y)$, which has no basic integral form. Rewriting it as $\sec^2(y)$ immediately unlocks the integral $\tan(y)$.
    \item \textbf{Exponentiating Logarithms to Isolate $y$ (Problem 9):} Moving from $\ln|y-4| = f(x) + C$ to $y - 4 = A e^{f(x)}$ by absorbing the $\pm e^C$ into a new constant $A$ is a cornerstone technique in solving ODEs explicitly.
\end{itemize}

\newpage
\section{Part 4: "Cheatsheet" and Tips for Success}

\subsection{Summary of Formulas}
\begin{itemize}
    \item \textbf{Separation Mechanism:} $\frac{dy}{dx} = \frac{f(x)}{g(y)} \implies g(y)dy = f(x)dx \implies \int g(y)dy = \int f(x)dx$
    \item \textbf{Logarithmic Integration:} $\int \frac{f'(x)}{f(x)} dx = \ln|f(x)| + C$
    \item \textbf{Integration by Parts:} $\int u \, dv = uv - \int v \, du$
\end{itemize}

\subsection{Cheats, Tricks, and Shortcuts}
\begin{itemize}
    \item \textbf{The $C$ Placement Shortcut:} You only need one $+C$ after integrating. By convention, append it to the side of the independent variable (usually $x$ or $t$). 
    \item \textbf{The Absolute Value Trick:} When integrating to a logarithm (e.g., $\ln|y|$), if you must solve for $y$, use the step $y = \pm e^C e^{f(x)}$ and immediately replace $\pm e^C$ with a new constant, say $A$. This absorbs the absolute value elegantly.
    \item \textbf{Solve for $C$ Early:} In IVPs, do not wait until you have explicitly solved for $y$ to plug in your points. Plug $(x_0, y_0)$ in \textit{immediately} after you perform the integration step. The algebra is almost always easier at this stage.
\end{itemize}

\subsection{Common Pitfalls and Mistakes}
\begin{itemize}
    \item \textbf{Fake Separation:} Attempting to separate by addition or subtraction (e.g., trying to say $\frac{dy}{dx} = x + y \implies dy - y = x \, dx$). This is mathematically illegal; separation \textit{must} be achieved purely via multiplication and division.
    \item \textbf{Forgetting the Constant Exponentiation:} A widespread error is writing $e^{x + C} = e^x + C$. The correct algebra is $e^{x + C} = e^x \cdot e^C = A e^x$.
    \item \textbf{Dropping the Negative in $u$-substitution:} In Problem 6, $u = -x^4 \implies du = -4x^3 dx$. Forgetting to divide by $-4$ destroys the rest of the problem.
\end{itemize}

\newpage
\section{Part 5: Conceptual Synthesis and The "Big Picture"}

\subsection{A) Thematic Connections}
The core theme of separable equations is \textbf{undoing the Chain Rule}. In Differential Calculus, when computing the derivative of a composite function like $\frac{d}{dx}[y^2]$, the chain rule gave us $2y \frac{dy}{dx}$. Separable differential equations reverse this exact process. By moving the $dy$ and $dx$ to opposite sides, we are isolating the ``inner'' derivative of a chain rule application, allowing us to perform standard Fundamental Theorem of Calculus integration independently on both variables.

\subsection{B) Forward and Backward Links}
\begin{itemize}
    \item \textbf{Backward Links:} This topic relies entirely on the algebraic robustness developed in Pre-Calculus (exponent rules, factoring) and the integration toolkit developed in Calculus II ($u$-sub, parts, partial fractions).
    \item \textbf{Forward Links:} Mastery here is non-negotiable for future chapters. Soon, you will encounter \textit{non-separable} linear first-order ODEs. You will learn the method of ``Integrating Factors'' (Section 2.5), which is essentially an algebraic trick to \textit{force} an inseparable equation to become separable. Later, in multivariable calculus and partial differential equations (PDEs), separation of variables is elevated to separating multivariable functions (like heat and wave equations) into products of single-variable functions.
\end{itemize}


\newpage
\section{Part 6: Real-World Application and Modeling}

\subsection{A) Concrete Scenarios in Finance and Economics}
Separable differential equations are the backbone of continuous-time financial modeling, economics, and stochastic calculus.
\begin{enumerate}
    \item \textbf{Continuous Wealth Accumulation (Compound Interest):} If wealth $W(t)$ grows at a continuous interest rate $r$, the rate of change of wealth is proportional to the current wealth: $\frac{dW}{dt} = rW$. This is the simplest separable equation and forms the foundation of all fixed-income mathematics.
    \item \textbf{Price Elasticity of Demand:} In microeconomics, elasticity $E$ measures how quantity demanded $Q$ changes with respect to price $P$. The formula is $E = \frac{P}{Q}\frac{dQ}{dP}$. If we assume a market with constant elasticity, we can rearrange this into a separable ODE to find the demand curve: $\frac{dQ}{dP} = E \frac{Q}{P}$.
    \item \textbf{The Vasicek Interest Rate Model (Deterministic Form):} In quantitative finance, interest rates don't just grow forever; they mean-revert. A simple deterministic version of an interest rate $r(t)$ reverting to a long-term mean $\mu$ at speed $a$ is $\frac{dr}{dt} = a(\mu - r)$. This is a separable equation used heavily in bond pricing models.
\end{enumerate}

\subsection{B) Model Problem Setup: Constant Elasticity Demand Curve}
\textbf{The Problem:} An economist determines that a particular luxury good has a constant price elasticity of demand of $-2$. Formulate a differential equation to find the demand function $Q(P)$. 
\newline
\textbf{Variables:}
\begin{itemize}
    \item $P = \text{Price of the good}$
    \item $Q = \text{Quantity demanded}$
\end{itemize}
\textbf{The Mathematical Model:}
Using the elasticity formula $E = \frac{P}{Q}\frac{dQ}{dP}$ and setting $E = -2$:
\[ -2 = \frac{P}{Q} \frac{dQ}{dP} \]
By separating variables, we set up the integral required to find the demand function:
\[ \int \frac{1}{Q} dQ = \int -\frac{2}{P} dP \]
Solving this ODE directly leads to the isoelastic demand curve function.

\newpage
\section{Part 7: Common Variations and Untested Concepts}

The homework set provided a solid foundation, but there are two incredibly important concepts regarding separable equations that were entirely absent from your assignment.

\subsection{Concept 1: Singular Solutions (The Lost Solution)}
When separating variables, we frequently divide both sides by a function of $y$, such as $p(y)$. However, in algebra, dividing by zero is undefined. If there exists a constant value $y = c$ such that $p(c) = 0$, dividing by $p(y)$ assumes $y \neq c$. By making this division, we might accidentally delete $y = c$ from our final answer. These lost constant solutions are called \textbf{singular solutions}.

\textbf{Example:} Solve $\frac{dy}{dx} = y^{2/3}$
\begin{itemize}
    \item \textit{Separation:} $y^{-2/3}dy = dx \implies 3y^{1/3} = x + C \implies y = \left( \frac{x+C}{3} \right)^3$.
    \item \textit{The Missing Piece:} Notice that early on, we divided by $y^{2/3}$. This is invalid if $y = 0$. If we check the original ODE, the constant function $y(x) = 0$ perfectly satisfies $\frac{dy}{dx} = 0 = (0)^{2/3}$. However, $y=0$ cannot be generated by our general solution formula for any value of $C$. Therefore, $y=0$ is a singular solution that must be stated separately!
\end{itemize}

\subsection{Concept 2: The Logistic Growth Equation (Partial Fractions)}
Many population and financial saturation models use the Logistic equation: $\frac{dy}{dt} = ky(1 - \frac{y}{L})$, where $L$ is a carrying capacity. This equation is separable but requires the method of \textbf{Partial Fractions} to integrate, which was not tested in your homework.

\textbf{Example Setup:} Solve $\frac{dy}{dt} = y(1 - y)$
\begin{itemize}
    \item \textit{Separation:} $\frac{1}{y(1-y)} dy = dt$.
    \item \textit{Integration Approach:} You cannot integrate the left side directly. You must decompose it using partial fractions:
    \[ \frac{1}{y(1-y)} = \frac{A}{y} + \frac{B}{1-y} \implies 1 = A(1-y) + By \implies A=1, B=1 \]
    \[ \int \left( \frac{1}{y} + \frac{1}{1-y} \right) dy = \int dt \implies \ln|y| - \ln|1-y| = t + C \]
    This leads to the classic S-curve shape fundamental to modeling market penetration and biological populations.
\end{itemize}

\newpage
\section{Part 8: Advanced Diagnostic Testing: "Find the Flaw"}

Below are 5 problems with full solutions. However, each solution contains a critical flaw. Your task is to play the role of the grader: identify the error, state why it is wrong in one sentence, and provide the correct calculation.

\subsection*{Problem A}
\textbf{Question:} Solve the differential equation $\frac{dy}{dx} = 5y$.
\newline
\textbf{Flawed Solution:}
\begin{align*}
    \frac{dy}{dx} &= 5y \\
    \frac{1}{y} dy &= 5 \, dx \\
    \int \frac{1}{y} dy &= \int 5 \, dx \\
    \ln|y| &= 5x \\
    y &= e^{5x} + C
\end{align*}

\subsection*{Problem B}
\textbf{Question:} Solve $\frac{dy}{dx} = \frac{x}{y^2}$.
\newline
\textbf{Flawed Solution:}
\begin{align*}
    \frac{dy}{dx} &= \frac{x}{y^2} \\
    \frac{1}{y^2} dy &= x \, dx \\
    \int \frac{1}{y^2} dy &= \int x \, dx \\
    \ln(y^2) &= \frac{1}{2}x^2 + C 
\end{align*}

\subsection*{Problem C}
\textbf{Question:} Determine if $\frac{dy}{dx} = x + xy$ is separable and separate it.
\newline
\textbf{Flawed Solution:}
\begin{align*}
    \frac{dy}{dx} &= x + xy \\
    dy &= (x + xy)dx \\
    dy - xy \, dx &= x \, dx \\
    \int dy - \int xy \, dx &= \int x \, dx
\end{align*}

\subsection*{Problem D}
\textbf{Question:} Solve $\frac{dy}{dx} = x y^3$.
\newline
\textbf{Flawed Solution:}
\begin{align*}
    \frac{dy}{dx} &= x y^3 \\
    \int y^3 dy &= \int x \, dx \\
    \frac{1}{4}y^4 &= \frac{1}{2}x^2 + C \\
    y^4 &= 2x^2 + C
\end{align*}

\subsection*{Problem E}
\textbf{Question:} Solve the IVP $\frac{dy}{dx} = -\frac{x}{y}, y(0) = -3$.
\newline
\textbf{Flawed Solution:}
\begin{align*}
    y \, dy &= -x \, dx \\
    \int y \, dy &= \int -x \, dx \\
    \frac{1}{2}y^2 &= -\frac{1}{2}x^2 + C \\
    y^2 &= -x^2 + C_1 \\
    (-3)^2 &= -(0)^2 + C_1 \implies 9 = C_1 \\
    y^2 &= -x^2 + 9 \\
    y &= \sqrt{9 - x^2}
\end{align*}

\vspace{0.5cm}
\hrule
\vspace{0.5cm}

\subsection*{Answer Key for "Find the Flaw"}

\textbf{Problem A Correction:}
\begin{itemize}
    \item \textbf{The Flaw:} The constant of integration $C$ was added \textit{after} exponentiating both sides, rather than immediately after integration.
    \item \textbf{Correct Step:} $\ln|y| = 5x + C \implies |y| = e^{5x+C} \implies y = A e^{5x}$.
\end{itemize}

\textbf{Problem B Correction:}
\begin{itemize}
    \item \textbf{The Flaw:} The student incorrectly integrated $y^{-2}$ as a natural logarithm instead of using the power rule.
    \item \textbf{Correct Step:} $\int y^{-2} dy = -y^{-1} = -\frac{1}{y}$. Thus, $-\frac{1}{y} = \frac{1}{2}x^2 + C$.
\end{itemize}

\textbf{Problem C Correction:}
\begin{itemize}
    \item \textbf{The Flaw:} The student attempted to separate variables using addition/subtraction, which is algebraically invalid for integration with respect to $dx$.
    \item \textbf{Correct Step:} Factor first. $\frac{dy}{dx} = x(1+y) \implies \frac{1}{1+y} dy = x \, dx$.
\end{itemize}

\textbf{Problem D Correction:}
\begin{itemize}
    \item \textbf{The Flaw:} The student multiplied by $y^3$ to move it to the left side, but algebraically they needed to divide by $y^3$.
    \item \textbf{Correct Step:} $\frac{1}{y^3} dy = x \, dx \implies \int y^{-3} dy = \int x \, dx \implies -\frac{1}{2}y^{-2} = \frac{1}{2}x^2 + C$.
\end{itemize}

\textbf{Problem E Correction:}
\begin{itemize}
    \item \textbf{The Flaw:} When taking the square root in the final step, the student assumed the positive root, contradicting the initial condition $y(0) = -3$.
    \item \textbf{Correct Step:} $y = \pm \sqrt{9 - x^2}$. Since $y(0)$ must be negative, the correct branch is $y = -\sqrt{9 - x^2}$.
\end{itemize}

\end{document}